\section{Methodology}
\label{sec:method}

\subsection{Overview}

Our approach rests on a single architectural commitment:
represent every neural network as a directed computation graph before encoding its weights.
This choice is motivated by the key insight ---
the graph schema (node=neuron, edge=weight) is architecture-agnostic in a way that flat
tokenization is not.
We combine this graph representation with permutation-equivariant GNN encoding and a
contrastive autoencoder SSL objective, trained on MLP/CNN model zoos and evaluated
zero-shot on ViT model zoos.

\subsection{Problem Setup}

Let $\mathcal{Z}_\text{train} = \{\theta_1, \ldots, \theta_N\}$ be a collection of trained
neural network checkpoints from architecture families $\mathcal{A}_\text{train}$
(MLPs and CNNs).
Let $\mathcal{Z}_\text{test} = \{\phi_1, \ldots, \phi_M\}$ be checkpoints from a held-out
architecture family $\mathcal{A}_\text{test}$ (ViT-S/16),
where $\mathcal{A}_\text{test} \cap \mathcal{A}_\text{train} = \emptyset$.
Our goal is to learn an encoder $f_\theta$ via SSL on $\mathcal{Z}_\text{train}$ such that
$f_\theta(\phi)$ is useful for predicting properties of held-out test models
$\phi \in \mathcal{Z}_\text{test}$ without any supervised signal from $\mathcal{Z}_\text{test}$.
We evaluate using a RidgeCV linear probe on extracted embeddings,
predicting ViT-S/16 ImageNet accuracy.

\subsection{Graph Encoding}

\textbf{Schema.}
We convert each checkpoint into a directed attributed graph $G = (V, E, \mathbf{x}_v, \mathbf{x}_e)$
where:
\begin{itemize}
  \item Nodes $v \in V$ correspond to neurons.
        Each node carries $\mathbf{x}_v = [\text{bias\_mean}, \text{bias\_std},
        \text{bias\_min}, \text{bias\_max}]$ (4-dim).
  \item Directed edges $(u, v) \in E$ represent scalar weights from neuron $u$ to neuron $v$.
        Each edge carries $\mathbf{x}_e = [\text{weight\_mean}, \text{weight\_std},
        \text{weight\_min}, \text{weight\_max}]$ (4-dim).
\end{itemize}

\textbf{Architecture homomorphism.}
All architecture families we consider ---
MLP linear layers, CNN convolutional filters, and ViT attention projections
(Q, K, V, MLP sublayers) ---
consist of linear maps between neuron sets.
A ViT attention head's weight matrix $\mathbf{W}_Q \in \mathbb{R}^{d_k \times d}$ defines
exactly the same edge structure as an MLP linear layer $\mathbf{W} \in \mathbb{R}^{d_\text{out}
\times d_\text{in}}$.
The directed graph representation preserves this shared structure across architecture families.

\subsection{Permutation-Equivariant SSL (EquiSSL-perm)}

\textbf{Encoder architecture.}
We use the neural-graphs encoder~\cite{Kofinas2024Graph} with 4 message-passing layers,
hidden dimension 256, and latent dimension 128.
This encoder is permutation-equivariant:
for any permutation $\pi$ of neurons within a layer,
$f_\theta(\pi(G)) = \pi(f_\theta(G))$.

\textbf{SSL objective.} We train with:
\begin{equation}
  \mathcal{L} = \mathcal{L}_\text{NT-Xent}(z, z^+) + \lambda \cdot \mathcal{L}_\text{MSE}(\hat{e}, e)
\end{equation}
where $z$ and $z^+$ are embeddings of two augmented views of the same network graph,
$\mathcal{L}_\text{NT-Xent}$ is the normalized temperature-scaled cross-entropy contrastive
loss~\cite{Chen2020Simple} with temperature $\tau = 0.07$,
and $\mathcal{L}_\text{MSE}$ is a reconstruction loss on edge attributes with $\lambda = 0.1$.

\textbf{Augmentation.}
Positive pairs are constructed by applying random neuron permutations to the same graph.
We do \emph{not} apply scale augmentation (see Section~\ref{sec:layernorm}).

\textbf{Training.}
Adam ($\text{lr} = 10^{-3}$, $\text{weight\_decay} = 10^{-4}$),
$\text{batch\_size} = 64$, 100 epochs,
CosineAnnealingLR ($T_\text{max} = 100$, $\eta_\text{min} = 10^{-5}$).

\subsection{Why Permutation-Only: The LayerNorm Gauge Hypothesis}
\label{sec:layernorm}

Our initial hypothesis followed Kalogeropoulos et al.~\cite{Kalogeropoulos2024Scale}
in using scale+permutation equivariance (the monomial matrix group).
The scale symmetry of a neural network is the equivalence $\theta \sim c\theta$ for scalar
$c > 0$ ---
two networks related by scaling of a neuron's incoming and outgoing weights are functionally
equivalent (for ReLU-like activations, which are scale-homogeneous).

We hypothesize that this symmetry breaks down in the presence of LayerNorm.
LayerNorm normalizes each neuron's activations to zero mean and unit variance at inference
time, which we argue fixes the scale of weight tensors.
Two ViT weight tensors that differ by scalar factor $c$ are \emph{not} functionally equivalent ---
LayerNorm normalizes them to the same activation scale regardless.
Therefore, scale equivariance would be an \emph{incorrect} inductive bias for ViT weight
encoders: it maps weight tensors that differ by scale to the same latent code,
discarding information that may distinguish functionally different ViT models.

This is the \textbf{LayerNorm gauge-fixing hypothesis} ---
an empirically motivated conjecture, not a proven mechanism.
The empirical evidence (EquiSSL underperforms EquiSSL-perm on ViT transfer;
see Section~\ref{sec:ablation}) is consistent with this hypothesis.
Direct validation would require comparing scale vs.\ permutation-only equivariance on
non-LayerNorm target architectures (e.g., ReLU MLP, BatchNorm CNN),
which we defer to future work.

Our primary method (EquiSSL-perm) uses permutation-only equivariance.
We include EquiSSL (scale+permutation, ScaleGMN backbone) as an ablation baseline.

\subsection{SANE Baseline with VICReg Fix}

As our primary baseline, we use SANE~\cite{Schurholt2024Towards} with a VICReg variance
regularization term~\cite{Bardes2021VICReg}:
\begin{equation}
  \mathcal{L}_\text{VICReg} = 0.1 \cdot \max(0,\; 1 - \sigma(z))
\end{equation}
Without this regularization, SANE's flat tokenizer maps all cross-architecture inputs to
near-constant latent codes ($\text{std} \approx 0.0014$).
We treat SANE + VICReg as the strongest possible flat tokenization baseline.

\subsection{Evaluation: RidgeCV Linear Probe}

After SSL training on $\mathcal{Z}_\text{train}$ (CNN checkpoints),
we extract embeddings $f_\theta(\phi)$ for all $\phi \in \mathcal{Z}_\text{test}$
(ViT checkpoints).
We fit a RidgeCV regressor ($\alpha \in \{0.1, 1.0, 10.0, 100.0\}$, 80/20 split)
to predict ViT-S/16 ImageNet accuracy.
Metric: $R^2$ (coefficient of determination).
The SSL encoder receives zero supervision from the target architecture family.
